\subsection{Allgemeines}
Die eingesetzten Reagenzien wurden von der zentralen Materialverwaltung bezogen. 
Das Kernresonanzspektrum wurde mit einem AV-250 der Fa. Bruker bei 300 K aufgenommen. \newline 
Die chemischen Verschiebungen sind in $\delta$-Werten (ppm) angegeben und beziehen sich für $^1H$-NMR-Spektren auf das Restprotonensignal des verwendeten Lösungsmittels \ce{CDCl3}. 
Bei der Zuordnung der Signale und für die Signalmultiplizitäten wurden die folgenden Abkürzungen verwendet: s – Singulett, bs – breites Singulett, d – Dublett, t – Triplett, q – Quartett, m – Multiplett.
Das IR-Spektrum wird mit einem Agilent Cary 630 FTIR gemessen. 

\subsection{Durchführung}
In einem Dreihalskolben wurden Styrol (\SI{5.21}{\gram}, \SI{50.0}{\milli\mole}), Chloroform (\SI{16.0}{\milli\liter}, \SI{200}{\milli\mole}), Benzyltriethylammoniumchlorid (\SI{0.114}{\gram}, \SI{0.50}{\milli\mole}) und Ethanol (\SI{0.5}{\milli\liter}) vorgelegt. 
Unter kräftigem Rühren wurde eiskalte, frisch bereitete 50\%ige Natronlauge (\SI{8.0}{\gram} NaOH in \SI{8.0}{\milli\liter} \ce{H2O}) zugegeben. 
Das zweiphasige Reaktionsgemisch wurde \SI{1}{\hour} bei Raumtemperatur und anschließend \SI{5}{\hour} bei \SI{50}{\degreeCelsius} gerührt.

Nach dem Abkühlen wurde das Gemisch in Wasser (\SI{250}{\milli\liter}) gegossen, die organische Phase abgetrennt und die wässrige Phase mit Chloroform (\SI{50}{\milli\liter}) extrahiert. 
Die vereinigten organischen Phasen wurden über wasserfreiem \ce{Na2SO4} getrocknet und filtriert. Es wurde mit weiteren \SI{50}{\milli\liter} Chloroform nachgespült. Das Lösungsmittel wurde am Rotationsverdampfer entfernt.
Das Rohprodukt wurde anschließend durch eine Vakuumdestillation gereinigt. 
Es wurde 1,1-Dichlor-2-phenylcyclopropan erhalten.